\chapter{More built-in types and the standard library}\label{ch-moretypes}

Chapters~\ref{ch-values} to~\ref{ch-procs} used four primitive types. This chapter introduces the
remaining built-in types of the language: complex numbers, numbers of
greater or smaller precision, bit sets and byte strings. It closes with a
survey of the constants and functions of the standard library. Much of
the chapter depends on the implementation, since the precision of numbers
does; where it does, the text says what ga68, a68j and Algol 68 Genie
provide.

\section{Complex numbers}\label{sec-complex}

The type \code|compl| holds a complex number, a pair of \code|real| values
for the real and the imaginary part. The binary operator \code|I| builds a
complex number from its two parts: \code|1 I 2| is $1 + 2i$. The
arithmetic operators \code|+|, \code|-|, \code|*| and \code|/| work on
complex numbers, also mixed with \code|int| and \code|real| operands, and
\code|=| and \code|/=| compare them. Further operators give the parts and
properties of a complex number:

\begin{center}
\begin{tabular}{ll}
\toprule
Operator & Result \\
\midrule
\code|RE z|, \code|IM z| & the real and the imaginary part \\
\code|ABS z| & the absolute value $|z|$ \\
\code|ARG z| & the argument (angle), between $-\pi$ and $\pi$ \\
\code|CONJ z| & the complex conjugate \\
\code|a I b| & the complex number $a + bi$ \\
\bottomrule
\end{tabular}
\end{center}

\example{ch06/complex}

The function \code|show| formats a complex number for printing. A
\code|real| (and therefore an \code|int|) is widened to \code|compl|
where a complex number is required, so \code|compl r := 2.5| is allowed;
this is the second step of the widening that Chapter~\ref{ch-values} introduced.

The standard library declares \code|compl| as a record,
\code|mode compl = struct (real re, im)|, so \code|re of z| selects the
real part of a complex variable and can be assigned to, and
\code|compl z = (1, 2)| is a record literal. Two details catch the
unwary: \code|I| has the highest precedence, 9, so \code|2 * x I y| means
\code|2 * (x I y)|; and complex numbers have no order, so there is no
\code|<|.

\begin{implnote}
ga68 does not yet provide the complex operators; a program that uses
them fails when it is linked. The examples in this section were run with
a68j and Algol 68 Genie. \code|Genie_Transput| does not print complex
numbers directly, which is why the example formats them itself; the
\code|print| of the Report and of Algol 68 Genie prints the two parts.
\end{implnote}

As an application, the following program solves quadratic equations,
whose roots may be complex. The discriminant decides whether the square
root is real or imaginary; the \code|if| balances its two branches to
\code|compl|.

\example{ch06/quadratic}

\section{Precision variants}\label{sec-long}

The types \code|int|, \code|real|, \code|compl| and \code|bits| come in
several sizes. The prefix \code|long| gives a variant with more bits,
\code|long long| one with still more, and \code|short| and
\code|short short| variants with fewer. How many variants there are and
how large they are depends on the implementation. A variant that the
implementation does not distinguish has the same values as the next
larger or smaller one, but it is still a type of its own: a
\code|long long int| must be converted with \code|SHORTEN| to become a
\code|long int|, even where both have 64 bits.

\begin{center}
\begin{tabular}{lccc}
\toprule
Type & ga68 & a68j & Algol 68 Genie \\
\midrule
\code|short short int| & 8 bits & 8 bits & 64 bits \\
\code|short int| & 16 bits & 16 bits & 64 bits \\
\code|int| & 32 bits & 32 bits & 64 bits \\
\code|long int| & 64 bits & 64 bits & 128 bits \\
\code|long long int| & 64 bits & 64 bits & 72 decimal digits \\
\midrule
\code|real| & 32 bits & 32 bits & 64 bits \\
\code|long real| & 64 bits & 64 bits & 128 bits \\
\code|long long real| & 80 bits & 64 bits & 63 decimal digits \\
\midrule
\code|bits| & 32 & 32 & 64 \\
\code|long bits| & 64 & 64 & 128 \\
\bottomrule
\end{tabular}
\end{center}

A literal of a long type is written with the prefix: \code|long 1|,
\code|long 2.5|, \code|long long 0|. The operator \code|LENG| converts a
value to the next longer type and \code|SHORTEN| to the next shorter one;
there is no implicit conversion between the sizes, so a \code|long int|
and an \code|int| cannot be added without \code|LENG|. The names of the
standard constants and functions for the long types begin with
\code|long_|: \code|long_max_int|, \code|long_pi|, \code|long_sqrt|.

\example{ch06/lengths}

The last two lines show the conversions: \code|LENG max_int + long 1|
computes in \code|long int| and so does not overflow, and \code|SHORTEN|
converts a value that fits back to \code|int|.

The long types are what a program needs for exact large integers and for
precise floating-point computations. With 32-bit integers, 13! already
overflows; with \code|long int|, all factorials up to 20! are exact:

\example{ch06/factorials}

\begin{implnote}
\code|Genie_Transput| prints \code|long int| and \code|long real| values,
but not the \code|short| and \code|long long| variants; they are
converted with \code|LENG| or \code|SHORTEN| before printing. In
Algol 68 Genie, \code|long long int| and \code|long long real| are
computed in software with a very large precision, which Rosetta Code
programs use for big numbers.
\end{implnote}

\section{Bit sets}\label{sec-bits}

A value of type \code|bits| is a fixed-size sequence of bits, as many as
the constant \code|bits_width| says. It is the Algol 68 equivalent of an
unsigned integer used as a set of flags or a bit mask. Literals are
written in base 2, 4, 8 or 16 with a prefix that gives the base:
\code|2r1100|, \code|8r14| and \code|16rc| are the same value. The
operators on \code|bits| are:

\begin{center}
\begin{tabular}{ll}
\toprule
Operator & Result \\
\midrule
\code|a AND b|, \code|a OR b|, \code|NOT a| & bitwise and, or, complement \\
\code|a SHL n|, \code|a SHR n| & shift left or right by \code|n| bits \\
\code|n ELEM a| & bit \code|n| of \code|a|, as a \code|bool| \\
\code|a = b|, \code|a /= b| & equality \\
\code|a <= b|, \code|a >= b| & subset, superset (not numerical order) \\
\code|ABS a|, \code|BIN n| & conversion to and from \code|int| \\
\bottomrule
\end{tabular}
\end{center}

\example{ch06/bits}

Two details differ from current conventions. First, \code|print| writes a
\code|bits| value as a row of \code|T| and \code|F| characters, one per
bit. Second, the bits are numbered from the left: bit 1 is the most
significant bit and bit \code|bits_width| the least significant one. In
\code|2r1100|, the two set bits are therefore numbers 29 and 30 of 32.
The function \code|bits_pack| builds a \code|bits| value from an array of
\code|bool|, whose last element becomes the least significant bit.

\begin{pitfall}
The comparisons of \code|bits| values are set comparisons:
\code|2r100 >= 2r011| is false, since the first value does not contain
the bits of the second. For numerical order, compare \code|ABS a| and
\code|ABS b|. \code|ELEM| has the precedence 7, so
\code|i + 1 ELEM b| means \code|i + (1 ELEM b)| and is rejected; write
\code|(i + 1) ELEM b|. A negative shift count shifts in the other
direction. And \code|ABS| of a value whose first bit is set does not fit
into an \code|int| of the same width: the Report leaves the result
undefined, and a68j gives a negative number.
\end{pitfall}

\begin{implnote}
ga68, a68j and Algol 68 Genie also provide \code|XOR| for \code|bits|,
which the Report does not have; it can be written as
\code|(a OR b) AND NOT (a AND b)|.
\end{implnote}

\section{Byte strings}\label{sec-bytes}

The type \code|bytes| is a string of fixed length, \code|bytes_width|
characters, stored compactly. It was meant for machines on which a
\code|string| was expensive. The function \code|bytes_pack| converts a
string to \code|bytes|, and \code|n ELEM b| selects a character.

\example{ch06/bytes}

\begin{implnote}
ga68 and a68j do not yet support \code|bytes|; the example was run with
Algol 68 Genie. Programs rarely need this type today, since \code|string|
serves the same purpose.
\end{implnote}

\section{Environment enquiries}\label{sec-enquiries}

The standard library contains constants that describe the implementation,
which the Report calls \rr{environment enquiries}. A program that wants to
be portable uses them instead of assuming particular sizes:

\example{ch06/enquiries}

The most important ones are:

\begin{center}
\begin{tabular}{ll}
\toprule
Constant & Meaning \\
\midrule
\code|max_int|, \code|long_max_int| & the largest integer of each size \\
\code|max_real|, \code|long_max_real| & the largest real of each size \\
\code|small_real| & the smallest \code|real| $x$ with $1 + x > 1$ and $1 - x < 1$ \\
\code|int_lengths|, \code|int_shorths| & 1 + the number of \code|long|, \code|short| sizes of \code|int| \\
\code|real_lengths|, \code|real_shorths| & the same for \code|real| \\
\code|bits_width|, \code|bytes_width| & the number of bits in \code|bits|, characters in \code|bytes| \\
\code|int_width|, \code|real_width| & the digits needed to print an \code|int|, a \code|real| \\
\code|max_abs_char| & the largest character code \\
\code|blank|, \code|null_character| & the space character, a filler character (code 0 here) \\
\code|flip|, \code|flop| & the characters for true and false (\code|"T"| and \code|"F"| here) \\
\bottomrule
\end{tabular}
\end{center}

\code|int_lengths| and \code|int_shorths| count \code|int| itself, so
\code|int_lengths = 2| means that there is one distinct \code|long| size
(and \code|int_shorths = 2| one \code|short| size); a larger prefix gives a type
with the same values as the largest distinct one.

\section{Mathematical functions}\label{sec-math}

The standard library provides the mathematical functions of
Section~\ref{sec-std-procs}, each also in a version for each long size:
\code|sqrt|, \code|exp|, \code|ln|, \code|sin|, \code|cos|, \code|tan|,
\code|arcsin|, \code|arccos|, \code|arctan|; for \code|long real|, they
are \code|long_sqrt|, \code|long_exp|, and so on. The function
\code|random| returns a pseudo-random \code|real| between 0 and 1.

\example{ch06/functions}

The negative widths ask \code|fixed| for a field of six or eight
characters without a plus sign (Section~\ref{sec-io}).

\section{Extending the standard operators}\label{sec-extend-ops}

A program can declare a standard operator for operand types that the
library does not cover, as Chapter~\ref{ch-procs} did for fractions. The standard
\code|MOD|, for example, is defined for integers only; a declaration adds
a version for \code|real| operands, while the standard \code|MOD| keeps its
meaning:

\example{ch06/real-mod}

\section*{Exercises}
\addcontentsline{toc}{section}{Exercises}

The \hyperref[ans-ch6]{answers} are at the end of the book.

\begin{enumerate}
\item What are \code|RE ((1 I 1) * (1 I -1))| and \code|ABS (0 I 2)|?
\item Why is \code|long int n := 5| not allowed, and how is it written
  correctly?
\item Assuming \code|bits_width| is 32, what are
  \code|ABS (16rff AND 2r1010)|, \code|ABS (2r1 SHL 4)| and
  \code|32 ELEM 2r1|?
\item Write a function \code|popcount| that counts the set bits in a
  \code|bits| value, using \code|ELEM|.
\item How can a program find out whether \code|long long int| is longer
  than \code|long int| in the implementation it runs on?
\item Is \code|2r100 >= 2r011| true? Why or why not?
\item What does \code|2 * 3 I 4| mean?
\end{enumerate}

\section*{Programming problems}
\addcontentsline{toc}{section}{Programming problems}

\begin{enumerate}
\item Compute $2^{62}$ and $3^{39}$ exactly with \code|long int|.
\item Print the first 90 Fibonacci numbers with \code|long int|.
\item Use the Newton iteration $z \leftarrow z - (z^3 - 1) / (3z^2)$
  to find which of the three cube roots of 1 is reached from each
  starting point of a grid in the complex plane, and print the grid as
  characters.
\item Represent a set of the numbers 1 to 32 as \code|bits|, and write the
  operations union, intersection, membership and size.
\item \textbf{Mandelbrot.} Print the Mandelbrot set with characters,
  using \code|compl| arithmetic.
\item \textbf{A bit sieve.} Write the sieve of Eratosthenes on an array
  of \code|bits| words, splitting a number into a word index and a bit
  number, and compare its size with the \code|[] bool| version of
  Chapter~\ref{ch-data}.
\item \textbf{Gray code.} Convert numbers to their Gray code
  (\code|n XOR (n SHR 1)| on \code|bits|) and back, and print the codes
  of 0 to 15.
\item \textbf{An adder.} Add two positive numbers given as \code|bits|
  bit by bit, as hardware does, with a carry, and report an overflow.
\item \textbf{Machine epsilon.} Find the smallest power of 2 that,
  added to 1, gives a result different from 1, by halving, for
  \code|real| and \code|long real|, and compare it with \code|small_real|.
\item \textbf{Exact factorials.} Compute 50! exactly with an array of
  \code|int| digits in base 10000, and compare the result with
  \code|long long int| in Algol 68 Genie.
\item \textbf{Safe fractions.} Make the rational numbers of
  Chapter~\ref{ch-procs} safe from overflow: compute the products in
  \code|long int| with \code|LENG|, reduce with a \code|long| greatest
  common divisor, and convert back with \code|SHORTEN| only after
  checking against \code|LENG max_int|. Use them for the first terms of a
  power series with exact coefficients, such as that of $e^x$.
\end{enumerate}
